% 49e Olympiade Internationale de Mathématiques — Madrid, Espagne.
% Énoncés officiels en français (imo-official.org). Chaque problème vaut 7 points.

\begin{center}
  \emph{Premier jour --- mercredi 16 juillet 2008 --- Durée : 4 heures 30}
\end{center}

\subsection*{Problème 1}
Soit $ABC$ un triangle dont les angles sont aigus, et soit $H$ son orthocentre. Le cercle passant par $H$ et
dont le centre est le milieu de $[BC]$ coupe la droite $(BC)$ en $A_1$ et $A_2$. De même, le cercle passant par
$H$ et dont le centre est le milieu de $[CA]$ coupe la droite $(CA)$ en $B_1$ et $B_2$, et le cercle passant
par $H$ et dont le centre est le milieu de $[AB]$ coupe la droite $(AB)$ en $C_1$ et $C_2$.

Montrer que $A_1$, $A_2$, $B_1$, $B_2$, $C_1$, $C_2$ sont cocycliques.

\subsection*{Problème 2}
(a) Montrer que
\[
\frac{x^2}{(x-1)^2} + \frac{y^2}{(y-1)^2} + \frac{z^2}{(z-1)^2} \geqslant 1
\]
pour tous nombres réels $x$, $y$, $z$, différents de $1$ et vérifiant $xyz = 1$.

(b) Montrer qu'il existe une infinité de triplets de nombres rationnels $x$, $y$, $z$, différents de $1$ et
vérifiant $xyz = 1$, pour lesquels l'inégalité ci-dessus est une égalité.

\subsection*{Problème 3}
Montrer qu'il existe une infinité d'entiers strictement positifs $n$ tels que $n^2 + 1$ possède un diviseur
premier strictement supérieur à $2n + \sqrt{2n}$.

\bigskip
\begin{center}
  \emph{Second jour --- jeudi 17 juillet 2008 --- Durée : 4 heures 30}
\end{center}

\subsection*{Problème 4}
Trouver toutes les fonctions $f$ de $]0, +\infty[$ dans $]0, +\infty[$ telles que
\[
\frac{\bigl(f(w)\bigr)^2 + \bigl(f(x)\bigr)^2}{f(y^2) + f(z^2)} = \frac{w^2 + x^2}{y^2 + z^2}
\]
pour tous nombres réels strictement positifs $w$, $x$, $y$, $z$, vérifiant $wx = yz$.

\subsection*{Problème 5}
Soient $n$ et $k$ des entiers strictement positifs tels que $k \geqslant n$ et $k - n$ est pair.

On suppose données $2n$ lampes numérotées de $1$ à $2n$ ; chacune peut être \emph{allumée} ou \emph{éteinte}.

Au début, toutes les lampes sont éteintes.

Une \emph{opération} consiste à allumer une lampe éteinte ou bien à éteindre une lampe allumée. On considère des
séquences constituées d'opérations successives.

Soit $N$ le nombre de séquences constituées de $k$ opérations et aboutissant à l'état où les lampes de $1$ à
$n$ sont allumées et les lampes de $n + 1$ à $2n$ sont éteintes.

Soit $M$ le nombre de séquences constituées de $k$ opérations et aboutissant à l'état où les lampes de $1$ à
$n$ sont allumées et les lampes de $n + 1$ à $2n$ sont éteintes, mais où les lampes de $n + 1$ à $2n$ n'ont
jamais été allumées.

Déterminer le rapport $N/M$.

\subsection*{Problème 6}
Soit $ABCD$ un quadrilatère convexe tel que $BA \neq BC$. Les cercles inscrits dans les triangles $ABC$ et
$ADC$ sont notés respectivement $\omega_1$ et $\omega_2$. On suppose qu'il existe un cercle $\omega$ qui est
tangent à la demi-droite $[BA]$ au-delà de $A$, tangent à la demi-droite $[BC]$ au-delà de $C$, et qui est
aussi tangent aux droites $(AD)$ et $(CD)$.

Montrer que les tangentes communes extérieures à $\omega_1$ et à $\omega_2$ se coupent en un point de $\omega$.
